\section{Convexity, noise scaling, and quantitative stability}
\label{sec:ext-stability}

\begin{proposition}[Convex sublevel sets and a genuine discontinuity]
\label{prop:ext-convexity}
For every $b_0\ge0$, the set $\{\rho:B_A^*(\rho)\le b_0\}$ is convex.
On each affine slice with a fixed Alice marginal, both $B_A^*$ and
$(B_A^*)^2$ are convex functions of $\rho$.
In general $B_A^*$ is neither convex nor continuous at states with
pure filtered-side marginal.
\end{proposition}

\begin{proof}
For a fixed invertible filter, filtering a mixture
$\sum_jq_j\rho_j$ produces the mixture of the individually filtered
states with weights proportional to their success probabilities.
The CHSH norm is convex, because it is the supremum of the appropriate
normalized linear Bell expressions. Its value on the filtered mixture
is therefore at most $\max_j B_A^*(\rho_j)$.
Taking the supremum over filters proves convexity of the sublevel set.

If all $\rho_j$ have the same Alice marginal, each fixed filter has
the same success probability on every component. The weights remain
$q_j$, and taking the supremum proves convexity of $B_A^*$ on the
slice. The square of a nonnegative convex function is convex,
giving the second assertion.

For the last statement, consider
\[
 \rho_p=p|\Phi^+\rangle\langle\Phi^+|
 +(1-p)|0\rangle\langle0|\otimes I/2.
\]
At $p=0$, Alice is pure and Bob maximally mixed, so $B_A^*(\rho_0)=0$.
For $0<p\le1/2$,
$\Delta_B(\rho_p)=p(2p-1)/2\le0$, which gives the upper bound one.
A rank-one Alice filter onto $|1\rangle$ removes the second component
and leaves Bob pure, giving the matching supremum one.
Thus $B_A^*(\rho_p)=1$ throughout $0<p\le1/2$.
It jumps at zero and exceeds the convex interpolation
$pB_A^*(|\Phi^+\rangle)=p\sqrt2$ for, for example, $p=1/4$.
\end{proof}

\begin{proposition}[Exact noise scaling on the unfiltered party]
\label{prop:ext-noise}
Let $\mathcal D_\eta(X)=\eta X+(1-\eta)\tr(X)I/2$ be the
depolarizing channel, with $0\le\eta\le1$. Then
\begin{equation}
 \boxed{
 B_A^*((\mathrm{id}\otimes\mathcal D_\eta)(\rho))
 =\eta B_A^*(\rho).
 }
 \label{eq:ext-noise-scaling}
\end{equation}
For $\eta>0$, the optimal Alice filters are unchanged.
If the original optimum exceeds one, the exact surviving-violation
threshold is $\eta>1/B_A^*(\rho)$.
For a canonical state the same statement applies to admixture of
global white noise.
\end{proposition}

\begin{proof}
Alice filtering commutes with the trace-preserving channel on Bob,
and the filtering success probability is unchanged.
For every normalized filtered state, the channel multiplies its
correlation matrix by $\eta$, and hence its CHSH norm by $\eta$.
The equality follows after taking suprema.
If $\rho_A=I/2$, then
$(\mathrm{id}\otimes\mathcal D_\eta)(\rho)
=\eta\rho+(1-\eta)I/4$, proving the last statement.
For a noncanonical state, global white noise does not generally
obey this scaling law.
\end{proof}

\begin{proposition}[A universal canonical Lipschitz estimate]
\label{prop:ext-lipschitz}
Let $S,S'$ be the squared unilateral optima of two canonical states,
with derived data $(A,s,d)$ and $(A',s',d')$. Then
\begin{equation}
 \boxed{
 |S-S'|\le
 2\norm{A-A'}_{\mathrm{op}}+|s-s'|+2\norm{d-d'}.
 }
 \label{eq:ext-lipschitz}
\end{equation}
This estimate requires no spectral gap and includes all singular
and repeated-eigenvalue cases.
\end{proposition}

\begin{proof}
For a common feasible pair, put $u=-n\times v$ and
$P=I-nn^T-vv^T$. Because $n\cdot v=0$, $\norm n=1$, and
$\norm v\le1$, the matrix $P$ is positive semidefinite and
$\tr P=2-\norm v^2\le2$. The objective is
\[
 f(n,v)=\tr(PA)+s\norm v^2+2u\cdot d.
\]
Since $\norm u=\norm v\le1$, the absolute difference of the
two objectives is bounded by the right-hand side of
\eqref{eq:ext-lipschitz}, uniformly over the common compact domain.
Taking maxima in both orders proves the claim.
\end{proof}

For example, an estimated value $\widehat S$ with certified right-hand
side at most $\varepsilon$ proves activation if
$\widehat S-\varepsilon>1$, and excludes activation if
$\widehat S+\varepsilon\le1$.
The estimate is in canonical coordinates; it does not assert uniform
continuity of whitening as a marginal approaches purity.

\begin{proposition}[Gradients at a simple optimum]
\label{prop:ext-gradient}
Suppose $H_0\succ0$ and the smallest eigenvalue of $M$ is simple.
Choose its normalized eigenvector $n$, put
$v=-K^{-1}Cn$, $u=-n\times v$, and
$P=I-nn^T-vv^T$. Then $S=(B_A^*)^2$ is differentiable in the
canonical data and
\begin{equation}
 \boxed{
 \nabla_D S=2(PD+uc^T),\qquad
 \nabla_c S=2(D^Tu+\norm v^2c).
 }
 \label{eq:ext-gradients}
\end{equation}
\end{proposition}

\begin{proof}
Under the stated hypotheses the optimizing reduced pair is unique up
to the simultaneous sign change $(n,v)\mapsto(-n,-v)$, which leaves
$P$ and $u$ unchanged. The simple-eigenvalue formula is smooth locally.
Equivalently, differentiating the maximized objective while holding
an optimizing pair fixed gives the same derivative.
In
$f=\tr(PDD^T)+\norm c^2\norm v^2+2u^TDc$,
differentiation with respect to $D$ and $c$ gives precisely
\eqref{eq:ext-gradients}.
\end{proof}

\paragraph{Independent checks and scope.}
The accompanying independent script tests these extensions through
complex $4\times4$ density-matrix calculations, separately from the
reference solver. With the recorded seed it checks the Gram identity,
plane-purity identity, determinant covariance, rational ellipsoid data,
and scalar probability formula on $1{,}000$ states of ranks one through
four. It tests $417$ purity-positive shortcut instances, $600$ physical
axial states, and $102$ success-frontier comparisons against direct
multistart filtering. The largest shortcut-attainment discrepancy is
$2.7\times10^{-15}$ and the largest frontier discrepancy is
$8.2\times10^{-10}$.
Further checks cover $1{,}000$ independent random rank-two states,
$129$ rational-cubic value evaluations, and $11$ points of the exact
tilted frontier. The rank-two ceiling discrepancy is below
$1.8\times10^{-15}$; the explicit tilted filters agree with direct
optimization within $2.9\times10^{-15}$.
These numerical checks support the mathematical
proofs; they do not establish them or extend the Lean coverage.
The exact formalization boundary is stated separately in
Appendix~\ref{sec:verification}.
